\documentclass[11pt,a4paper]{article}
% Shared preamble for RuPhO English problem statements (match T1 2026 style).
% Use: \input{latex/rupho_en_preamble.tex} after \documentclass.
\usepackage[margin=1in]{geometry}
\usepackage{amsmath,amssymb}
\usepackage{graphicx}
\usepackage{float}
\usepackage{enumitem}
\usepackage{microtype}
\usepackage{caption}
\usepackage{tabularx}
\usepackage{hyperref}

\captionsetup{font=small,labelfont=bf,skip=6pt}

\setlist[itemize]{leftmargin=1.2cm}
\setlist[enumerate]{leftmargin=1.2cm}

% One outer box per subproblem: [id | statement | points] (no inner rules)
\newcommand{\problembox}[3]{%
  \par\addvspace{0.42em}%
  \noindent
  \setlength{\fboxsep}{7.5pt}%
  \setlength{\fboxrule}{0.95pt}%
  \fbox{%
    \begin{minipage}{\dimexpr\linewidth-2\fboxsep-2\fboxrule\relax}
    \begin{tabularx}{\linewidth}{@{}>{\raggedright\arraybackslash\bfseries}p{2.1em} @{\hspace{0.9em}} >{\raggedright\arraybackslash}X @{\hspace{0.9em}} >{\raggedleft\arraybackslash\bfseries}p{2.4em}@{}}
    #1 & #2 & #3
    \end{tabularx}
    \end{minipage}%
  }%
  \par\addvspace{0.32em}%
}


\title{\textbf{Discrete model of avalanche motion}}
\author{\normalsize Y12 (2012) --- English translation}
\date{}
\begin{document}
\maketitle

On a long inclined plane with angle α, heavy cubes are placed at regular intervals $L$. They are held by the force of static friction.

\begin{figure}[H]
    \centering
    \includegraphics[width=0.72\textwidth]{T2_EN_assets/fig1.jpg}
    \caption{Figure 1.}
\end{figure}

The model assumes that if the cube is given a negligible speed, its further movement occurs without friction. If the top cube starts to move, it occurs without friction. If the top cube moves, it will collide with the second cube, then the chain of two cubes will collide with the third, and so on. All collisions are assumed to be absolutely inelastic. As a result, a long chain of cubes appears, to which more and more cubes are added. This process models the movement of an avalanche along a mountain slope.

\problembox{A1}{Let $n = 100$ cubes move in the chain. Determine the increment in the speed $\Delta v_1$ of the chain immediately after the collision with the $101$th cube compared to the speed after the collision with the $100$th cube.}{0}

\problembox{A2}{Find the difference in speed $\Delta v_2$ of chains of $100$ and $400$ cubes.}{0}

\problembox{A3}{What changes in the movement of an avalanche when taking into account the friction force? Assuming that slopes with angles $\alpha \geq 23^\circ$ are avalanche dangerous, give answers to the previous two questions taking into account the friction force. At this point, the slope angle is considered equal to $\alpha = 23^\circ$.}{0}

\vspace{1em}
\noindent\textit{Source:} \url{https://pho.rs/p/3977}

\end{document}